\documentclass[10pt,a4paper]{article}
\usepackage[T1]{fontenc}
\usepackage[margin=2cm]{geometry}
\usepackage{isabelle,isabellesym}
\usepackage{amsmath,amssymb}
\usepackage{pdfsetup}

\urlstyle{rm}
\isabellestyle{it}

\begin{document}

\title{Conway's Game of Life: Locality, Gliders and the Gosper Gun}
\author{Carles Mar\'in}
\maketitle

\begin{abstract}
We formalize Conway's Game of Life~\cite{Gardner70,BCG82} on the integer plane: a generation is an arbitrary set
of live cells, the neighbourhood is the Moore neighbourhood and the rule is B3/S23. We prove that the step commutes
with every bijection of the plane that preserves the neighbourhood relation (translations, the quarter turn and the
diagonal reflection, hence the symmetry group of the square), the light-cone locality of the dynamics (after $t$
generations the state on a region $R$ depends only on the initial state on $R$ thickened by $t$ cells), that finite
patterns stay finite, and that patterns whose $t$-thickenings are disjoint evolve independently for $t$ generations.
The last result is extended to arbitrary families of patterns, which gives the free flight of salvos: translated
copies of a spaceship, in arbitrary phases, that stay apart during one period fly independently for ever (for
gliders it suffices that the offsets differ by at least 5 in some coordinate). We prove Conway's speed limit: a
pattern contained in a diagonal half-plane needs two generations to advance one diagonal, so a finite non-empty
pattern that reappears displaced by $(a,b)$ after $p$ generations satisfies $2(|a|+|b|)\le p$; the glider and the
lightweight spaceship attain the diagonal and the orthogonal bounds.
Concrete evolutions are proved by trace replay: an external simulator proposes the intermediate generations and a
verified executable step, evaluated by \texttt{code\_simp}, checks every transition. We prove that the block is a
still life, that the blinker has period~2, that the glider and the lightweight spaceship have period~4 with
displacements $(1,1)$ and $(-2,0)$, that eater~1 absorbs a glider in ten generations, and that the Gosper glider gun
emits a glider every 30 generations: after $30k$ generations the gun is intact and is followed by $k$ gliders in
flight, so its population is $36+5k$. Hence there is a finite pattern whose population is unbounded, and grows at
least linearly.
\end{abstract}

\paragraph{Acknowledgement.} The formalization and the text of this document were prepared with the
assistance of generative AI systems, and checked by the author.

\tableofcontents

\section*{Introduction}

The theory \texttt{Life\_Base} defines the rule on sets of cells of $\mathbb{Z}^2$ and proves equivariance,
locality (the cone lemma) and finiteness; it also defines a simple executable step on lists of cells, proved
equal to the abstract one, and the trace-replay theorem used for every concrete evolution.
\texttt{Life\_Objects} proves the block, the blinker and the glider (period, minimal period, free flight for every
number of periods). \texttt{Life\_Disjoint} proves that patterns whose thickenings stay disjoint evolve
independently, the formulation in terms of regions that contain the orbits, and an executable sufficient
condition. \texttt{Life\_Flight} extends the one-generation lemma to arbitrary families of patterns and proves the
free flight of salvos of spaceships. \texttt{Life\_Speed} proves the speed limit~\cite[Theorem~4.1]{JG22}: the
half-plane lemmas hold for arbitrary patterns, finite or not, and the bound for a finite pattern that reappears
displaced is stated in the oblique form of~\cite[Exercise~4.26]{JG22}; nothing is stated for objects that never
reappear exactly. \texttt{Life\_Fast} gives a faster executable step on sorted lists (merging the eight translated copies
and counting by run-length encoding); it checks at run time that its input is sorted and otherwise falls back to the
simple step, so its correctness holds for every list. \texttt{Life\_Spaceships} proves the lightweight spaceship (which
attains the orthogonal bound of the speed limit) and eater~1. \texttt{Gosper\_Gun} proves the gun theorem and the unbounded population.

The theories \texttt{Life\_Spaceships\_Data} and \texttt{Gosper\_Gun\_Data} contain only data: the generations
proposed by the external simulator. They are not trusted; \texttt{Gosper\_Gun\_Cert} and the concrete lemmas check
every transition. No oracle is used.

The disjointness lemma for a single generation is the one stated by Myreen and Carneiro (their area of
influence is our thickening by one cell) in their verified construction, in HOL4, of the Game of Life inside the
Game of Life~\cite{MC25}; here it is iterated for $t$ generations. The patterns are the classical
ones; see~\cite{JG22} for their history and properties. Coordinates: $x$ is the column, growing to the right, and
$y$ the row, growing downwards.

\parindent 0pt\parskip 0.5ex

\input{session}

\bibliographystyle{abbrv}
\bibliography{root}

\end{document}
